Autonomous lunar landing is a control problem with direct operational consequences: during the powered-descent phase it requires precise positioning and safe contact, in real time, over cratered, rocky, and sloped terrain. In recent years GPU-parallel physics engines (NVIDIA PhysX) and the simulation frameworks built on them have made it practical to learn complex, contact-rich tasks with deep reinforcement learning (RL) by running thousands of environments concurrently \citep{mittal2023,nvidia2025,rudin2022}. The project underlying this study belongs to that class: it is an Isaac Lab direct-workflow task in which a four-legged rocket with a two-axis gimbal-thrust engine learns, using the Soft Actor-Critic (SAC) algorithm, to land on a procedural lunar surface built from a real digital elevation model (DEM) macro surface enriched with craters, rocks, and centimetre-scale regolith texture.

In procedural multi-resolution terrain of this kind, two layers are typically used to keep the compute and memory budget bounded: a low-resolution global surface covering a wide area, and a high-resolution local patch generated only in the region near the agent or the target where it is needed. Seamlessly joining these two layers is a classic problem, studied in the computer-graphics literature for more than twenty-five years (geometry clipmaps: \citealp{losasso2004}; ROAM: \citealp{duchaineau1997}; geomipmapping: \citealp{deboer2000}). That literature, however, focuses almost entirely on rendering quality; how a discontinuity at the join boundary affects foot-to-ground contact, altimeter readings, or the noise character of an RL agent's proprioceptive observation is rarely addressed quantitatively.

In the system studied here this problem appears with two distinct and mutually inconsistent solutions: the training environment blends between layers with a compact-support cubic \texttt{smoothstep} function, whereas a previously used mesh-based implementation uses an infinite-support Gaussian bell. These two functions differ in a fundamental mathematical property: \texttt{smoothstep} reaches exactly zero at a finite radius (compact support), whereas the standard Gaussian bell never reaches exactly zero at any finite radius (infinite support; \citealp{hardy1971}) and must be truncated somewhere to be usable in an engineering application; that truncation risks creating a new, measurable seam at the cut point. Comparing the two methods fairly requires forcing both to zero at the same transition radius (hard truncation) and measuring them under the same transition budget; this form of comparison is not directly addressed in any of the terrain-randomization studies surveyed here (see Section~\ref{sec:related}).

Motivated by these observations, this study addresses three research questions:

\begin{itemize}
  \item RQ1: When compared fairly under the same transition-radius budget, how large is the quantitative difference in boundary continuity (C0/C1) between compact-support (\texttt{smoothstep}) and infinite-support (Gaussian) blending functions, and what trade-off does this difference create with computational and memory cost?
  \item RQ2: Can a \texttt{hybrid} blending method that uses a Gaussian in the landing-gear contact region and \texttt{smoothstep} beyond it combine the continuity advantages of both methods at a reasonable increase in computational cost?
  \item RQ3: Do the continuity differences among these three blending methods translate into a measurable difference in learning performance in a real Isaac Sim/Isaac Lab SAC training run?
\end{itemize}

Accordingly, the aim of the study is threefold: (i) to compare the existing \texttt{smoothstep} method and the previously used Gaussian method quantitatively within the same architecture and under the same transition budget; (ii) to propose and validate a \texttt{hybrid} method that combines the advantages of both; and (iii) to compare these three methods not only at the analytic/CPU level but inside a real Isaac Sim/Isaac Lab GPU-parallel physics simulation, with a real SAC training run, and thereby show whether the choice of blending function is reflected in learning performance. This last point is central to the novelty of the study: the comparison is carried out not as an abstract numerical experiment but on real lunar terrain, inside a real Isaac Sim training loop.

The main contributions of this study are as follows:

\begin{itemize}
  \item The existing \texttt{smoothstep} and the previously used Gaussian blending methods are, for the first time, compared fairly under the same transition budget by truncating them at a finite transition radius (hard truncation); it is shown that the untruncated Gaussian appears well-behaved but in fact carries a hidden seam defect that emerges once it is truncated at a finite radius.
  \item A new \texttt{hybrid} blending method is proposed that uses a switch weighted toward a Gaussian in the contact region and toward \texttt{smoothstep} beyond it; this method is shown to provide continuity at machine precision and to create no new seam at its own internal transition.
  \item A statistically validated (Wilcoxon signed-rank test) comparison of continuity, cost, and VRAM is performed over thirty seeds and four method variants (\texttt{smoothstep}, as-shipped Gaussian, calibrated Gaussian, \texttt{hybrid}); the robustness of the findings is tested with four separate ablation analyses (contact radius, switch width, transition radius, grid resolution).
  \item The three methods are compared inside a real Isaac Lab/Isaac Sim PhysX environment, with GPU-parallel SAC training; unlike prior work in the terrain-randomization literature, this is the first attempt to treat the choice of blending function directly as a training-design variable.
  \item This real training comparison is run not with a single seed per method but with three independent seeds (n=3); by showing that the large apparent performance gap seen in the initial single-seed measurement cannot be confirmed statistically with multiple seeds, the study provides a concrete, self-sourced cautionary example about the reliability of single-run RL comparisons.
  \item All three methods are integrated into the production code in a backward-compatible way (without changing existing behaviour) and are selectable via an environment variable (\texttt{ISAACLAB\_TERRAIN\_BLEND\_METHOD}).
\end{itemize}

The remainder of this paper is organized as follows. Section~\ref{sec:related} reviews the literature in which the study is situated, under five thematic clusters. Section~\ref{sec:method} details the experimental environment and the simulation setup shared by all runs, the mathematical definition of the blending functions, the analytic evaluation protocol, the reinforcement learning method used (the task's MDP formulation and SAC), and the real Isaac Sim training protocol. Section~\ref{sec:results} presents the continuity, cost, VRAM, ablation, and real training results. Section~\ref{sec:discussion} discusses the findings in relation to the literature and addresses the limitations. Section~\ref{sec:conclusion} summarizes the study and closes with concrete directions for future work.
